\documentclass[10pt,a4paper]{amsart}
\usepackage[margin=19mm]{geometry}
\usepackage{amsmath,amssymb,mathtools}
\usepackage[scr=rsfs,cal=euler]{mathalfa}
\usepackage{xfrac}
\usepackage[unicode,hidelinks,bookmarks=false]{hyperref}
\newcommand{\ie}{i.e.\ }
\newcommand{\cf}{cf.\ }
\newcommand{\resp}{respectively\ }
\newcommand{\Subsection}{\subsection}
\providecommand{\nobreakdash}{\nobreak\hbox{-}}
\setlength{\emergencystretch}{2em}
\multlinegap0pt
\usepackage{amssymb}
\usepackage{xcolor}
\usepackage{tikz}
\newtheorem{lemma}{Lemma}
\title{Permutation representations on cohomology of toric varieties}
\author{Tao Gui, Chushi Qin, Kaizhe Shen,, Rui Xiong}
\date{Source: arXiv:2609.06597v1 (CC BY 4.0)}
\begin{document}
\maketitle

\noindent\emph{Source and licence.} Permutation representations on cohomology of toric varieties. Version arXiv:2609.06597v1 (CC BY 4.0), distributed by arXiv under the Creative Commons Attribution 4.0 International licence, http://creativecommons.org/licenses/by/4.0/, as recorded in the arXiv licence metadata for this exact version and shown on the abstract page https://arxiv.org/abs/2609.06597v1. Full licence notice: \texttt{licenses/arxiv-2609.06597-CC-BY-4.0.txt}.\par\smallskip

\noindent\textit{Editorial scope.} Counted unit: the article's lemma lem:orbit-fixed (source0.tex lines 632--639). The article's complete original proof of it is delivered with it (source0.tex lines 641--673, one \texttt{proof} environment of 33 lines). No mathematics is written, repaired or re-derived: the statement and the proof are the article's own lines, in the article's own notation, and no retained line is substituted or re-flowed.  No further source-local context is needed: every cross-reference of every retained line resolves inside the two delivered spans.  Nothing was omitted from any retained span, so the package carries no omission record.  No retained line contains a citation, so the wrapper carries no bibliography and no bibliography entry.  No retained line contains a figure environment, an image-inclusion command or drawing code, so no image is delivered and the package declares no asset dependency.  Editorial formatting only, in addition: an AMS \texttt{amsart} preamble that restates the article's own declarations and macros used by the retained lines, namely the environments \texttt{lemma} on separate counters, together with the standard packages those lines need.  The retained environments print in this document as Lemma 1 in order of appearance, whereas the article prints the same items with the numbering of its own full document.  The complete native source of the article as distributed in the arXiv e-print for this exact version is bundled unchanged as \texttt{sources/209552/source0.tex}.
\begin{lemma}\label{lem:orbit-fixed}
For every $\sigma\in\Sigma$ such that $g\sigma=\sigma$,
\begin{equation*}
  \chi_c\bigl(O(\sigma)^g\bigr)
  =\det\bigl(I-\psi_\sigma(g)\bigr).
\end{equation*}
No properness assumption on the action is needed.
\end{lemma}

\begin{proof}
Put
\[
N_\sigma=N\cap\operatorname{span}_\mathbb{R}(\sigma),
  \qquad L_\sigma=N/N_\sigma.
\]
The subgroup $N_\sigma$ is saturated, and $O(\sigma)$ is the torus with
cocharacter lattice $L_\sigma$.  The element $g$ induces a lattice
automorphism $A$ of $L_\sigma$, whose realification is
$\psi_\sigma(g)$.  The fixed locus $O(\sigma)^g$ is the kernel of the torus
homomorphism induced by $A-I$.

Apply the Smith normal form to $I-A$.  If its rank is $s$ and its nonzero
invariant factors are $d_1,\ldots,d_s$, then, non-canonically,
\begin{equation*}
  O(\sigma)^g
  \simeq
  \mu_{d_1}\times\cdots\times\mu_{d_s}
  \times(\mathbb{C}^\times)^{\operatorname{rank}(L_\sigma)-s}.
\end{equation*}
If $I-A$ is singular, the positive-dimensional torus factor makes
$\chi_c$ vanish.  If $I-A$ is nonsingular, the fixed locus has
$|\det(I-A)|$ points.

Since $g$ has finite order, every real eigenvalue of $A$ is $1$ or $-1$,
and the non-real eigenvalues occur in conjugate pairs on the unit circle.  A
$1$-eigenvalue makes the determinant zero, a $-1$-eigenvalue contributes
$2$, and a pair $\lambda,\overline\lambda$ contributes
\[
  (1-\lambda)(1-\overline\lambda)=|1-\lambda|^2>0.
\]
Thus $\det(I-A)\geq0$, and the absolute value may be removed.
\end{proof}

\end{document}
